% !TeX root = ../paper.tex
\section{Representation Probe Diagnosis Protocol}
\label{appendix:representation-diagnostics}

\subsection{Probe Setup and Paired Interventions}

The probes test whether failures of instruction following coexist with accessible target information
in the action-trained representations. If this information is absent, strengthening its
correspondence with actions alone cannot exploit an existing semantic representation. Conversely,
reliable decoding motivates supervision that links the available information to the required
behavior. The probes measure accessibility in the final frozen policies, rather than changes from
their original VLMs or the causal contribution of an individual action head.

We freeze $\pi_{0.5}$, GR00T N1.6, and StarVLA and train a common-form
conditional target-point probe over each model's multimodal fusion
representation. The probe receives language and visual representations
associated with the current command and predicts a two-dimensional target
point. Behavior and probe results use the same paired instances.

Sem switches fine-grained commands while holding image and robot state
fixed, testing whether the target representation changes with semantics.
Vis changes scene layout or moves an explicit visual cue, testing whether
the representation binds to the current visual position. Con switches
target commands in an action-contrast subset, testing whether the internal
target representation updates with the control requirement. Sem uses
fine-grained commands from the original task templates, Vis uses explicit
visual targets under a unified protocol, and Con uses paired commands in
the action-contrast subset.

All three models use matched instance pairs, target annotations, and probe-training
budgets; the feature location is selected at the corresponding multimodal
fusion layer.

\subsection{Representation Metrics}

Let $\hat{\mathbf{x}}$ be the predicted target point, $\mathbf{x}$ the
current target, and $\mathbf{x}'$ the alternative candidate, expressed in
the same evaluation-image coordinate system. The per-instance target-point
error is $e=\lVert\hat{\mathbf{x}}-\mathbf{x}\rVert_2$ in pixels. Target
assignment measures the proportion of predictions satisfying
$\lVert\hat{\mathbf{x}}-\mathbf{x}\rVert_2 <
\lVert\hat{\mathbf{x}}-\mathbf{x}'\rVert_2$; its two-choice reference is
50\%, and it is distinct from behavioral object selection or task completion.
Representation response direction measures agreement between the change
in predicted target position and the true target displacement across a
paired language or visual intervention.

These metrics measure the readability of target information in frozen
representations and its response to interventions. They do not directly
measure whether actions follow that information and are therefore
interpreted jointly with the behavioral results in Appendix~\ref{behavioral-diagnosis}.

\subsection{Results and Interpretation}
\label{representation-diagnosis}

\Needspace{0.28\textheight}
% !TeX root = ../paper.tex
\begin{center}
\setlength{\belowcaptionskip}{\baselineskip}
\captionof{table}{Behavioral and representation diagnostics for $\pi_{0.5}$, GR00T
N1.6, and StarVLA. Obj., Pos., and Comp. are behavioral rates (\%), Err. is
target-point error in pixels, Assign. is target assignment (\%), and Resp. is
representation response direction.}
\label{tab:diagnostics}
\normalsize
\setlength{\tabcolsep}{4pt}
\renewcommand{\arraystretch}{1.10}
\begin{tabular}{@{}llrrrrrr@{}}
\toprule
Model & Cond. & Obj.$\uparrow$ & Pos.$\uparrow$ & Comp.$\uparrow$ &
Err.$\downarrow$ & Assign.$\uparrow$ & Resp.$\uparrow$ \\
\midrule
$\pi_{0.5}$ & ORI & 100.00 & 100.00 & 99.00 & -- & -- & -- \\
$\pi_{0.5}$ & Sem & 88.00 & 11.00 & 6.00 & 14.19 & 96.63 & 0.927 \\
$\pi_{0.5}$ & Vis & 83.00 & 14.00 & 5.00 & 20.26 & 91.09 & 0.873 \\
$\pi_{0.5}$ & Con & -- & 25.00 & 25.00 & 12.02 & 97.50 & 0.978 \\
GR00T N1.6 & ORI & 100.00 & 100.00 & 98.00 & -- & -- & -- \\
GR00T N1.6 & Sem & 79.00 & 5.00 & 1.00 & 12.26 & 96.72 & 0.967 \\
GR00T N1.6 & Vis & 80.00 & 3.00 & 2.00 & 18.35 & 92.11 & 0.896 \\
GR00T N1.6 & Con & -- & 13.00 & 11.00 & 16.18 & 93.89 & 0.912 \\
StarVLA & ORI & 100.00 & 100.00 & 96.00 & -- & -- & -- \\
StarVLA & Sem & 86.00 & 3.00 & 2.00 & 11.93 & 98.82 & 0.981 \\
StarVLA & Vis & 72.00 & 9.00 & 4.00 & 23.14 & 83.80 & 0.832 \\
StarVLA & Con & -- & 6.00 & 6.00 & 19.66 & 89.64 & 0.877 \\
\bottomrule
\end{tabular}
\end{center}
%

Table~\ref{tab:diagnostics} reports the paired behavioral and representation diagnostics. Target
assignment rates are 96.63--98.82\% under Sem, 83.80--92.11\% under Vis, and 89.64--97.50\% under
Con, all above the 50\% two-choice reference. The corresponding response-direction scores are
0.927--0.981, 0.832--0.896, and 0.877--0.978. These results indicate that the frozen representations
retain decodable target information that varies with fine-grained commands, explicit visual prompts,
and updated control requirements.

However, decodability does not imply that the policy effectively uses this information to generate
actions. The low behavioral completion rates, despite positive target-assignment and
response-direction results, reveal a gap between target information in the representations and
execution that follows the current instruction. This contrast supports action-aligned supervision
connecting current instructions and visual prompts to local actions, but does not by itself
establish complete instruction understanding or identify a unique cause of execution failure.

